\subsection{The finite XX certificate with a shared nine-dimensional register}
\label{app:xx-assisted-certificate}

Let \(L_N\) denote the loss of the explicit strategy on
\(\psi_{\mathrm{XX},N}\otimes\Phi_9\). Repeating each Schmidt amplitude nine
times and dividing it by \(3\) commutes with duplication by two and preserves
the norm. The ninefold blocks in Theorem~\ref{thm:operational-upper} therefore
give \(L_N=2\cE(\psi_{\mathrm{XX},N})/243\). We bound this loss at
\(N=2^{18},2^{19},2^{20}\).

\paragraph{Covariance compression.}
Use the exact parity kernel~\eqref{eq:xx84-kernel}, whose singular list
occurs twice, and the dyadic quadrature above with order
\(k=\min(24,L)\) on a bin of length \(L\).
The polynomial comparison leading to~\eqref{eq:xx84-compression} now
uses degree \(23\). For \(m=N/4\le2^{18}\), the reciprocal kernel
\(1/(\pi w)\) has Frobenius norm at most
\(\sqrt{1+\ln(2m)}/(2\pi)<1\) on the integer grid and less than \(2\)
on the quadrature grid. The same geometric and Cauchy remainder estimates
with this order give
\[
 \|s(B)-s(K)\|_2
 \le 3\,3^{-24}
      +\tfrac12\bigl[(3/2)3^{-24}+(8/3)2^{-24}\bigr]
 <7.95\times10^{-8}.
\]
Here \(K\) is the exact weighted quadrature kernel and singular lists are
padded by zeroes. The dimensions of \(K\) are \(296,320,344\), respectively.

\paragraph{Arithmetic and spectral errors.}
The verifier \texttt{assisted\_xx\_certificate.py} uses floating-point
eigensolvers to propose roots and eigenpairs. On each bin with \(L>24\)
it evaluates the rational Jacobi determinant polynomials by
\(p_0(x)=1,\ p_1(x)=x,\ p_{j+1}(x)=xp_j(x)-\beta_j^2 p_{j-1}(x)\).
The 24 disjoint sign-changing brackets for \(p_{24}\) are bisected on
the grid \(10^{-28}\mathbb Z\). The resulting brackets enclose every
Jacobi eigenvalue. For a root \(\mu\), the Jacobi eigenvector with first
entry one has entries \(p_j(\mu)/(\beta_1\cdots\beta_j)\). The quadrature
weight is \(L\) divided by the sum of their squares. This gives intervals
for the exact nodes and weights.

With directed arithmetic on \(10^{-45}\mathbb Z\), the code encloses
\(\pi\) by Machin's formula and evaluates the degree-24 cosine and sinc
polynomials on the argument intervals. For \(|t|\le2\) the alternating
remainders are bounded by \(2^{26}/26!<10^{-18}\). The entry enclosures
give a symmetric midpoint matrix \(\widehat K\) and
\(\|K-\widehat K\|_{\mathrm F}<3.32\times10^{-18}\) in all three runs.

For each of 24 proposed eigenvectors \(v\), the verifier takes \(\mu\)
to be the nearest grid point to its exact Rayleigh quotient. Integer
arithmetic gives an upper bound on \(\|\widehat K v-\mu v\|_2/\|v\|_2\).
The disjoint residual intervals contain distinct eigenvalues.
Subtracting the squares of the lower bounds on their absolute values
from the exact \(\|\widehat K\|_{\mathrm F}^2\) bounds the omitted squared
singular mass by \(3\times10^{-15}\). For a model ratio \(r_i=q^{k_i}\)
with \(q=32767/32768\), the model singular value is \(s_i^*=r_i/(1+r_i^2)\).
Table~\ref{tab:xx-assisted-exponents} gives the three lists of 24 exponents.
The verifier encloses their ratios by directed binary exponentiation. The sum of squared differences
between the matched model entries and the retained singular values is
less than \(1.21\times10^{-11}\). Sorting the matched lists, with omitted
entries paired with zero, can only decrease their \(\ell^2\) distance.
Combining the three errors gives the one-parity distance
\(\|s(B)-s^*\|_2<3.55\times10^{-6}\) in all three runs. The verifier uses
the individual enclosures for each size.

The physical and model singular values lie below \(1/2\). A singular value
\(s\) gives the mode Schmidt angle \(\theta(s)=\tfrac12\arcsin(2s)\).
On \([0,s_{\max}]\) the derivative is at most
\((1-4s_{\max}^2)^{-1/2}\), with \(s_{\max}\) bounded by the largest
model value plus the one-parity error. Placing the model in the physical
Schmidt mode bases gives squared vector distance
\(2(1-\prod_i\cos^2(\theta(s_i)-\theta(s_i^*)))\), which is at most
\(2\sum_i(\theta(s_i)-\theta(s_i^*))^2\). This uses both copies of the
one-parity list. Rearrangement again decreases the distance. The verifier
checks that the resulting bound on the ordered amplitude distance fits
the fixed rational budget \(d\) in Table~\ref{tab:xx-assisted-certificate}.

\begin{table}[htb]
\centering\footnotesize
\begin{tabular}{@{}lp{13cm}@{}}
\toprule
$N$ & Exponents $k_i$ in increasing order \\
\midrule
$2^{18}$ & 11037, 33596, 57136, 81569, 106688, 132338, 158423, 184879, 211661, 238735, 266074, 293658, 321469, 349490, 377711, 406118, 434702, 463456, 492369, 521437, 550651, 580008, 609500, 639124 \\
$2^{19}$ & 10537, 32018, 54344, 77460, 101191, 125398, 149994, 174920, 200135, 225608, 251316, 277239, 303362, 329670, 356152, 382798, 409599, 436548, 463637, 490860, 518212, 545687, 573281, 600990 \\
$2^{20}$ & 10080, 30585, 51822, 73761, 96253, 119175, 142448, 166016, 189843, 213900, 238165, 262622, 287256, 312054, 337007, 362104, 387338, 412703, 438192, 463798, 489518, 515347, 541280, 567314 \\
\bottomrule
\end{tabular}
\caption{The three finite geometric models for the assisted certificates. Each exponent occurs twice; the common amplitude ratio is $32767/32768$.}
\label{tab:xx-assisted-exponents}
\end{table}

\paragraph{Exact ranks and score transfer.}
Each ratio \(q^{k_i}\) occurs twice in the normalized geometric model.
The amplitudes are \(\sqrt{p_0}q^e\), where
\(p_0=\prod_i(1+q^{2k_i})^{-2}\), in increasing order of \(e\), with
multiplicities from \(\prod_i(1+z^{k_i})^2\).
Keep the unnormalized prefix \(u\) at \(e\le400000\). The integer
recurrence and rank interval intersections used for the \(2^{84}\)
model give its mass and duplication overlap in directed arithmetic.
Put \(E_0=\|(I-D)u\|_2^2\) and \(\tau=1-\|u\|_2^2\).
The full model's overlap with its duplication is at least the prefix
overlap by positivity, so its duplication defect is at most
\(E_0+2\tau\). Since \(D\) is an isometry, the two norm comparisons
with the full model give
\[
 \max\{0,\sqrt{E_0}-2(d+\sqrt{\tau})\}
 \le \sqrt{\cE(\psi_{\mathrm{XX},N})}
 \le \sqrt{E_0+2\tau}+2d.
\]
The upper bound thus pays the omitted geometric mass linearly before
transferring the square root of the defect. The verifier checks that each
omitted prefix mass satisfies \(\tau<1.30\times10^{-8}\) and uses the fixed
\(d\) budget in each row to obtain the loss intervals.

\begin{table}[htb]
\centering\small
\begin{tabular}{@{}lrrll@{}}
\toprule
\(N\) & Prefix rank & \(d\) upper bound & Certified interval for \(\cE(\psi_N)\) &
  Certified interval for \(L_N\)\\
\midrule
\(2^{18}\) & \(4375\) & \(1.363\times10^{-5}\) &
  \([0.0906646,\,0.0908053]\) &
  \([0.000746210786,\,0.000747368515]\)\\
\(2^{19}\) & \(5642\) & \(9.187\times10^{-6}\) &
  \([0.0852473,\,0.0853883]\) &
  \([0.000701624233,\,0.000702783850]\)\\
\(2^{20}\) & \(7295\) & \(1.684\times10^{-5}\) &
  \([0.0803123,\,0.0804798]\) &
  \([0.000661007046,\,0.000662385080]\)\\
\bottomrule
\end{tabular}
\caption{Certified embezzlement errors of the exact XX ground states
$\psi_N$ and losses of the prescribed strategy on $\psi_N\otimes\Phi_9$.
The intervals round the verifier's endpoints outwards. The amplitude-distance
bound $d$ includes covariance compression, spectral truncation and arithmetic
errors.}
\label{tab:xx-assisted-certificate}
\end{table}

The floor \(\gamma=(3-2\sqrt2)/243\) lies between
\(0.0007060612150362\) and \(0.0007060612150363\).
The certified assisted losses are below \(0.000703\) and \(0.000663\)
at \(N=2^{19}\) and \(N=2^{20}\), respectively.
